\documentclass{ltc-guji}
\setmainfont{FandolSong}
\关闭分页

\begin{document}
\begin{正文}

% ============================================================================
% 工尺谱 Level 1 — 主字 + 右侧工尺字串 (issue #90 step A)
% ============================================================================

% 基本：1 主字 + 1 工尺音
\工尺{春}{尺}单音

% 1 主字 + 2 工尺音
\工尺{眠}{六五}两音

% 1 主字 + 3 工尺音
\工尺{晓}{尺工六}三音

% 1 主字 + 4 工尺音
\工尺{鸟}{尺工六五}四音

% 句中混合 — 多个工尺连用
\工尺{春}{尺}\工尺{眠}{工}\工尺{不}{六}\工尺{觉}{五}\工尺{晓}{尺工}春眠

% ============================================================================
% 自定义参数
% ============================================================================

% 缩放 — 工尺字更大
\工尺[缩放=0.6]{大}{尺工}大尺工

% 缩放 — 工尺字更小
\工尺[缩放=0.4]{小}{尺工六}小三音

% 颜色
\工尺[颜色=red]{红}{尺工}红色工尺

\工尺[颜色=blue, 缩放=0.55]{兰}{六五}蓝色工尺

% 偏移 — 工尺字离主字更远 (向内偏负值)
\工尺[偏移=-2pt]{远}{尺工}远偏

% ============================================================================
% 全局设置
% ============================================================================

\工尺设置{颜色=red, 缩放=0.55}
\工尺{东}{尺}全局红
\工尺{风}{工六}全局红
\工尺设置{颜色=black, 缩放=0.5}

% ============================================================================
% 与正文交错 — 一句歌词 + 工尺
% ============================================================================

\工尺{春}{尺}\工尺{眠}{工}\工尺{不}{六}\工尺{觉}{五}\工尺{晓}{尺工}%
\工尺{处}{六}\工尺{处}{五}\工尺{闻}{尺}\工尺{啼}{工}\工尺{鸟}{六五}

% ============================================================================
% Level 2 — \音 命令配 板眼 装饰 (#90 step B/C/D)
% ============================================================================

% 单 \音 不带板眼 — 等同于纯字符
\工尺{春}{\音{尺}}单音

% 单 \音[板] — 工尺字 + 红色板点
\工尺{眠}{\音[板]{尺}}单板

% 单 \音[眼] — 红色眼点
\工尺{不}{\音[眼]{工}}单眼

% 多个工尺，每个独立板眼
\工尺{觉}{\音[板]{尺}\音[眼]{工}}板+眼

% 多个工尺，混合 板/眼/无
\工尺{晓}{\音[板]{尺}\音{工}\音[眼]{六}}三音三态

% 头板/底板/底叮
\工尺{春}{\音[头板]{尺}}头板
\工尺{眠}{\音[底板]{尺}}底板
\工尺{不}{\音[底叮]{尺}}底叮

% 板眼颜色覆盖
\工尺{觉}{\音[板, 颜色=blue]{尺}}蓝板

% 完全手动 (不用预设)
\工尺{晓}{\音[字符=●, 缩放=0.25, 颜色=green]{尺}}手动绿点

% 完整唐诗 + 板眼 — 体现工尺谱真实排版
\工尺{春}{\音[板]{尺}\音[眼]{工}}%
\工尺{眠}{\音[板]{六}}%
\工尺{不}{\音[眼]{五}\音[板]{尺}}%
\工尺{觉}{\音[眼]{工}}%
\工尺{晓}{\音[板]{六}\音[眼]{五}}春眠不觉晓

% ============================================================================
% Level 3 — 八度标记 (#90 Level 3)
% ============================================================================

% 单 \音[高] — 仅八度，无板眼
\工尺{高}{\音[高]{尺}}单高

% 单 \音[低]
\工尺{低}{\音[低]{尺}}单低

% 单 \音[昆低]
\工尺{昆}{\音[昆低]{尺}}昆低

% 板 + 高 (双槽位同时使用)
\工尺{合}{\音[板, 高]{尺}}板+高

% 眼 + 低
\工尺{合}{\音[眼, 低]{工}}眼+低

% 多个工尺，每个独立板眼+八度
\工尺{月}{\音[板, 高]{尺}\音[眼, 低]{工}}多音双槽

% 八度颜色覆盖
\工尺{蓝}{\音[高, 八度颜色=blue]{尺}}蓝高

% 全自定义八度 (无预设)
\工尺{绿}{\音[八度字符=亻, 八度颜色=green]{尺}}绿亻

% ============================================================================
% Level 3 续 — 拖音 (延音) 测试
% ============================================================================

% 单 拖音 — 一拍延音 (无板眼)
\工尺{延}{\音{尺}\拖音}1拍延

% 拖音带板眼 — 每延一拍带节拍标记
\工尺{长}{\音[板]{尺}\拖音[眼]\拖音[板]\拖音[眼]}4拍尺

% 多拖音 + 八度
\工尺{高}{\音[板, 高]{尺}\拖音[眼]\拖音[板]}3拍高尺

% 完整短句：每个字 1-4 拍
\工尺{春}{\音[板]{尺}\拖音[眼]}%
\工尺{眠}{\音[板]{工}\拖音[眼]\拖音[板]\拖音[眼]}%
\工尺{晓}{\音[板]{六}\拖音[眼]}1+4+2拍

\end{正文}
\end{document}
